\begin{abstract}
Public access to a satellite archive does not guarantee enough eligible observations in a selected benchmark. We assess a fixed benchmark of coastal ultra-high-resolution Advanced Scatterometer observations from the National Oceanic and Atmospheric Administration, using a workflow that recorded no access to reference buoy winds. The design comprised 128 satellite files and 12 stations, with eight evaluation stations assessed over five consecutive 28-day blocks. Readiness required at least 400 unique station--observation opportunities overall and at least eight in every station--block cell. After processing 96 primary files and four prespecified rounds of eight additional files, the design yielded 150 opportunities. The minimum cell count was zero, and only three of 40 cells reached eight. Wind-performance evaluation therefore did not proceed. The original repository record reports zero reads of National Data Buoy Center wind outcomes; a later deterministic replay found no outcome paths or fields and reproduced the capacity decision, terminal counts, transfer accounting, and predictor-file integrity. In the 1,024 evaluation station--file records, 628 ended without a spatial match under the specified rule. The five-stage acquisition sequence used 9.418 decimal GB. These results describe the selected files and processing rules, rather than archive-wide coverage or physical causes of missing observations. Because the audit was formulated after the capacity failure, its findings are exploratory. The study provides an inspectable readiness procedure and a negative case in which a benchmark was stopped at its predictor-only stage. Passing the subsequent integrity audit did not change the failed readiness decision.
\end{abstract}
